\subsection{Cold Plate and Water Path}
\label{sec:cooling}

Each submodule sits on one liquid-cooled cold plate. The heat capacity of the plate and the water keeps the heat sink
temperature constant over one period, so the cooling path is calculated with the mean losses of the period. The
$n_{par}$ modules in parallel at a switch position have parallel water paths with the flow $Q_w / n_{par}$ each,
$Q_w$ the flow per cold plate. Along one path the water passes the module of one switch position and then the module
of the other (T1/D1 first in the example), so the second module sees water warmed by the first.

The thermal resistance heat sink to water under one module is one value of the converter design, $R_{th\_hw\_ref}$
at the reference flow $Q_{ref}$, with the flow dependence of the heat transfer coefficient in turbulent flow~[3]:
\begin{equation}
  R_{th\_hw}(Q_w) = R_{th\_hw\_ref} \cdot \left( \frac{Q_{ref}}{Q_w} \right)^{0.8} .
\end{equation}
The example uses $R_{th\_hw\_ref} = 6$\,K/kW at $Q_{ref} = 10$\,l/min, a value to be replaced by the data of the
cold plate. For module $j$ along the path with the mean losses $P_j$ (IGBT and diode),
\begin{align}
  \Delta T_{w\_j} &= \frac{P_j}{\rho \cdot c_p \cdot Q_w / n_{par}} , &
  T_{w\_j} &= T_{w\_in} + \sum_{i<j} \Delta T_{w\_i} + \frac{\Delta T_{w\_j}}{2} , \\
  T_{hs\_j} &= T_{w\_j} + P_j \cdot R_{th\_hw}(Q_w) , &
  T_{case\_j} &= T_{hs\_j} + P_j \cdot R_{th\_cs} ,
\end{align}
with $\rho \cdot c_p = 4.13$\,kJ/(l$\cdot$K) for water at about 50\,$^{\circ}$C ($988\,\mathrm{kg/m^3} \cdot
4.18\,\mathrm{kJ/(kg \cdot K)}$) and $T_{w\_j}$ the mean water temperature along the module. The water inlet
temperature $T_{w\_in}$ (default 50\,$^{\circ}$C) and the flow per cold plate $Q_w$ (default $Q_{ref}$) are part of the
operating point. The water outlet temperature follows from the losses of the cold plate,
$T_{w\_out} - T_{w\_in} = P_{plate} / (\rho \cdot c_p \cdot Q_w)$.
